\section{Phase-averaged form factor and classical symmetry defects}
\label{sec:interaction-phase-form-factor}
The phase average below is an average over an auxiliary \(16\)-torus of kick
phases. It is not a theorem about the fixed kick, and it does not bear on
the signed off-diagonal sum. The canonical proof is
\path{mathematics/fixed-map/interaction_phase_form_factor.md}.
The rational identities are certified by
\path{formalization/Bgs/PhaseFormFactorRational.lean} and recomputed by the
Python and C++ verifiers named there. No external refereeing is asserted.
BGS-TRACE-005 and BGS-BGS-001 remain open.

\begin{theorem}[Phase-averaged form factor, BGS-TRACE-126]
\label{thm:interaction-phase-form-factor}
Let \(N\ge 3\) be odd, let \(\eta\in\mathbb R\), and let each of the \(16\)
pairs of nonzero Fourier modes of the kick carry an independent uniform phase.
For every integer \(m\ge 1\),
\[
\mathbb E_\varphi\lvert\operatorname{Tr} U_N(\eta,\varphi)^m\rvert^2
=\lvert K_m\rvert+\eta^2 D_2(N,m)+\rho_m(N,\eta),
\]
where \(K_m=\ker(S^m-I)\) and \(D_2(N,m)\) is the explicit Gauss-sum
coefficient of the proof. It satisfies
\(\lvert D_2\rvert\le 2m^2\sigma\lvert K_m\rvert\) with
\(\sigma=66033/160000\). If \(x=m B\lvert\eta\rvert\) and
\(B=184981/160000\), the averaged remainder satisfies
\[
\lvert\rho_m(N,\eta)\rvert
\le N^2\Biggl(\frac{x^3}{6}\Bigl(2+\frac{x^3}{6}\Bigr)+\frac{x^4}{4}\Biggr).
\]
The fixed-kick remainder through the same order is at most
\(N^2\bigl[(x^3/6)(2+x^3/6)+x^3+x^4/4\bigr]\).

At an odd prime \(p\), with \(m\) the order of \(S\) modulo \(p\),
\[
D_2(p,m)=-m\bigl(\sigma-\sigma_0(p)\bigr)p^2,
\]
where \(\sigma_0(p)\) is the energy of the symbol modes congruent to \(0\)
modulo \(p\). This energy equals \(81/64000\) at \(p=3\), \(1/320000\) at
\(p=5\), \(81/200\) at \(p=7\), and \(0\) for every prime \(p\ge 11\).
Thus \(D_2(3,4)=-1184949/80000\), \(D_2(5,6)=-79239/1280\),
\(D_2(7,3)=-181251/160000\) and \(D_2(41,5)=-111001473/32000\).
At the single phase \(\varphi=0\) the linear
coefficient need not vanish: at \(N=11\), \(m=3\) it is nonzero by a
cyclotomic witness.

Parity \(\psi(x)\mapsto\psi(-x)\) commutes with \(M_N\), and position-basis
complex conjugation reverses \(M_N\). At \(N=7\) the kick equals \((9/10)I\),
so both symmetries survive for every real coupling. For every other odd
\(N\) both defects are positive. For \(0<\lvert\eta\rvert<\pi/B\) the kicked
propagator therefore fails both symmetries outside dimension \(7\). The
parity defect equals \(663/1000000\) once \(N\ge 19\). The reversal defect
equals \(66031/80000\) once \(N\ge 79\).

No revival limit is proved. A positive limit of \(\tau^2 N^3\) at times
comparable to \(N\) lies outside the truncation \(m B\lvert\eta\rvert\le 1\).
The Kurlberg--Rudnick moment calculation is not an input. An \(N\)-dependent
antiunitary other than the two symmetries above is not classified.
\end{theorem}
